% Relatorio em LaTeX para inclusao no artigo principal
% Uso sugerido no artigo: % Relatorio em LaTeX para inclusao no artigo principal
% Uso sugerido no artigo: % Relatorio em LaTeX para inclusao no artigo principal
% Uso sugerido no artigo: % Relatorio em LaTeX para inclusao no artigo principal
% Uso sugerido no artigo: \input{relatorio_reprodutibilidade}

\section{Independent Reproducibility Report of the RCTM-based PRNG}
\label{sec:repro_report}

This section summarizes an independent reproducibility assessment of the Robust Chaotic Tent Map (RCTM) random number generator implementation available in this repository. The objective is to verify whether the observed statistical behavior is consistent across independent execution, and to document both positive and negative findings with full transparency.

\subsection{Repository and Data Availability}

All source code, test scripts, raw test results, and documentation are publicly available at:

\texttt{https://github.com/angualberto/rctm-reproducibility-study}

The repository includes complete test batteries, binary data files (\texttt{agle\_v2\_10m.bin}, \texttt{rctm61.81\_c.bin}), Dieharder results, LaTeX source for this report, and comprehensive documentation (README, quickstart guide, research findings). This enables full independent verification and replication of all findings presented in this study.

\subsection{Experimental Setup}

The experiments were executed on Linux x86\_64 using IEEE-754 double precision arithmetic and C/Python tooling available in the repository. The main analyzed stream was \texttt{agle\_v2\_10m.bin} (10 MB), with additional statistics extracted from:

\begin{itemize}
\item Dieharder full battery (\texttt{-a}): \texttt{agle\_v2\_dieharder\_rctm.txt}
\item Python statistical scripts: \texttt{test\_*.py}
\item Consolidated local report: \texttt{RESULTADOS\_TESTES\_PYTHON\_2026-03-06.txt}
\end{itemize}

\subsection{Technique and Evaluation Protocol Used}

The technical protocol combined first-order distribution checks, sequence-level dependency tests, and distinguishability analysis. The adopted workflow was:

\begin{enumerate}
\item \textbf{Data generation and ingestion}: binary output stream generated by the RCTM implementation and parsed as \texttt{uint8} and \texttt{uint32} depending on each test.
\item \textbf{First-order randomness checks}: Shannon entropy and bit-balance verification to quantify marginal distribution quality.
\item \textbf{Order-pattern complexity}: permutation entropy (order 3) to assess ordinal diversity in local windows.
\item \textbf{Predictability assessment}: N-gram next-symbol prediction (2-gram, 3-gram, 4-gram) with comparison against random baseline (1/256 = 0.391\%).
\item \textbf{Distinguishing capability}: pattern-window entropy test to evaluate whether the generator can be statistically distinguished from an ideal random source.
\item \textbf{Dependency structure}: LAG-1 pair-space coverage and linear autocorrelation across lags 1, 2, 4, 8, 16, and 32.
\item \textbf{State-space probe}: nearest-neighbor state-reconstruction indicator in embedded space.
\item \textbf{External battery validation}: Dieharder \texttt{-a} full suite for broad-spectrum stress testing.
\end{enumerate}

The interpretation criteria followed practical thresholds used in the local scripts: for N-gram, values substantially above baseline indicate exploitable dependence; for autocorrelation, values near zero across multiple lags indicate absence of strong linear dependence; for LAG-1 coverage, low occupancy of the 256$^2$ pair space indicates constrained transitions.

\subsection{Summary of Main Results}

Table~\ref{tab:repro_summary} presents the key outcomes obtained in the independent rerun.

\begin{table}[htbp]
\centering
\caption{Independent reproducibility summary for RCTM output.}
\label{tab:repro_summary}
\begin{tabular}{lcc}
\hline
Test / Indicator & Measured value & Assessment \\
\hline
Shannon entropy & 7.998454 bits/byte & PASS \\
Permutation entropy (order 3) & 0.9999998568 & PASS \\
N-gram predictability (2-gram) & 1.630\% & FAIL \\
Distinguishing attack (normalized entropy) & 0.5190 & FAIL \\
LAG-1 pair coverage & 78.46\% & FAIL \\
Bit balance (ones ratio) & 49.99\% & PASS \\
Autocorrelation (lags 1..32) & near 0 & PASS \\
State reconstruction indicator & no recoverable state & PASS \\
\hline
Overall (Python battery) & 6/9 passed & MIXED \\
\hline
\end{tabular}
\end{table}

In Dieharder, multiple failures/weak results were observed, including \texttt{diehard\_opso}, \texttt{diehard\_squeeze}, repeated failures in \texttt{marsaglia\_tsang\_gcd}, multiple failures in \texttt{rgb\_lagged\_sum}, and failures in \texttt{dab\_bytedistrib}, \texttt{dab\_filltree2}, and \texttt{dab\_monobit2}.

\subsection{Interpretation}

The independent reproduction indicates a \textbf{local-versus-structural randomness gap}. Specifically:

\begin{itemize}
\item \textbf{Local quality appears strong}: high Shannon entropy, good bit balance, and low linear autocorrelation.
\item \textbf{Sequence-level structure remains detectable}: failing N-gram predictability, low pair-space coverage (LAG-1), and distinguishability from ideal random behavior.
\end{itemize}

This combination suggests that, although single-byte statistics are close to ideal, there are non-linear or higher-order dependencies across consecutive outputs.

\subsection{Reproducibility Statement}

All scripts, logs, and result files used to obtain these numbers are included in the repository, enabling direct independent verification. No filtering of unfavorable outcomes was applied in this report.

\subsection{Conclusion for the Reproducibility Study}

From the independent execution, the evidence supports the following conclusion: the evaluated RCTM stream shows good first-order statistical properties but does not consistently satisfy stronger unpredictability/indistinguishability expectations under extended batteries. Therefore, the current implementation should be treated with caution for cryptographic use until the identified sequence-level issues are addressed and revalidated.


\section{Independent Reproducibility Report of the RCTM-based PRNG}
\label{sec:repro_report}

This section summarizes an independent reproducibility assessment of the Robust Chaotic Tent Map (RCTM) random number generator implementation available in this repository. The objective is to verify whether the observed statistical behavior is consistent across independent execution, and to document both positive and negative findings with full transparency.

\subsection{Repository and Data Availability}

All source code, test scripts, raw test results, and documentation are publicly available at:

\texttt{https://github.com/angualberto/rctm-reproducibility-study}

The repository includes complete test batteries, binary data files (\texttt{agle\_v2\_10m.bin}, \texttt{rctm61.81\_c.bin}), Dieharder results, LaTeX source for this report, and comprehensive documentation (README, quickstart guide, research findings). This enables full independent verification and replication of all findings presented in this study.

\subsection{Experimental Setup}

The experiments were executed on Linux x86\_64 using IEEE-754 double precision arithmetic and C/Python tooling available in the repository. The main analyzed stream was \texttt{agle\_v2\_10m.bin} (10 MB), with additional statistics extracted from:

\begin{itemize}
\item Dieharder full battery (\texttt{-a}): \texttt{agle\_v2\_dieharder\_rctm.txt}
\item Python statistical scripts: \texttt{test\_*.py}
\item Consolidated local report: \texttt{RESULTADOS\_TESTES\_PYTHON\_2026-03-06.txt}
\end{itemize}

\subsection{Technique and Evaluation Protocol Used}

The technical protocol combined first-order distribution checks, sequence-level dependency tests, and distinguishability analysis. The adopted workflow was:

\begin{enumerate}
\item \textbf{Data generation and ingestion}: binary output stream generated by the RCTM implementation and parsed as \texttt{uint8} and \texttt{uint32} depending on each test.
\item \textbf{First-order randomness checks}: Shannon entropy and bit-balance verification to quantify marginal distribution quality.
\item \textbf{Order-pattern complexity}: permutation entropy (order 3) to assess ordinal diversity in local windows.
\item \textbf{Predictability assessment}: N-gram next-symbol prediction (2-gram, 3-gram, 4-gram) with comparison against random baseline (1/256 = 0.391\%).
\item \textbf{Distinguishing capability}: pattern-window entropy test to evaluate whether the generator can be statistically distinguished from an ideal random source.
\item \textbf{Dependency structure}: LAG-1 pair-space coverage and linear autocorrelation across lags 1, 2, 4, 8, 16, and 32.
\item \textbf{State-space probe}: nearest-neighbor state-reconstruction indicator in embedded space.
\item \textbf{External battery validation}: Dieharder \texttt{-a} full suite for broad-spectrum stress testing.
\end{enumerate}

The interpretation criteria followed practical thresholds used in the local scripts: for N-gram, values substantially above baseline indicate exploitable dependence; for autocorrelation, values near zero across multiple lags indicate absence of strong linear dependence; for LAG-1 coverage, low occupancy of the 256$^2$ pair space indicates constrained transitions.

\subsection{Summary of Main Results}

Table~\ref{tab:repro_summary} presents the key outcomes obtained in the independent rerun.

\begin{table}[htbp]
\centering
\caption{Independent reproducibility summary for RCTM output.}
\label{tab:repro_summary}
\begin{tabular}{lcc}
\hline
Test / Indicator & Measured value & Assessment \\
\hline
Shannon entropy & 7.998454 bits/byte & PASS \\
Permutation entropy (order 3) & 0.9999998568 & PASS \\
N-gram predictability (2-gram) & 1.630\% & FAIL \\
Distinguishing attack (normalized entropy) & 0.5190 & FAIL \\
LAG-1 pair coverage & 78.46\% & FAIL \\
Bit balance (ones ratio) & 49.99\% & PASS \\
Autocorrelation (lags 1..32) & near 0 & PASS \\
State reconstruction indicator & no recoverable state & PASS \\
\hline
Overall (Python battery) & 6/9 passed & MIXED \\
\hline
\end{tabular}
\end{table}

In Dieharder, multiple failures/weak results were observed, including \texttt{diehard\_opso}, \texttt{diehard\_squeeze}, repeated failures in \texttt{marsaglia\_tsang\_gcd}, multiple failures in \texttt{rgb\_lagged\_sum}, and failures in \texttt{dab\_bytedistrib}, \texttt{dab\_filltree2}, and \texttt{dab\_monobit2}.

\subsection{Interpretation}

The independent reproduction indicates a \textbf{local-versus-structural randomness gap}. Specifically:

\begin{itemize}
\item \textbf{Local quality appears strong}: high Shannon entropy, good bit balance, and low linear autocorrelation.
\item \textbf{Sequence-level structure remains detectable}: failing N-gram predictability, low pair-space coverage (LAG-1), and distinguishability from ideal random behavior.
\end{itemize}

This combination suggests that, although single-byte statistics are close to ideal, there are non-linear or higher-order dependencies across consecutive outputs.

\subsection{Reproducibility Statement}

All scripts, logs, and result files used to obtain these numbers are included in the repository, enabling direct independent verification. No filtering of unfavorable outcomes was applied in this report.

\subsection{Conclusion for the Reproducibility Study}

From the independent execution, the evidence supports the following conclusion: the evaluated RCTM stream shows good first-order statistical properties but does not consistently satisfy stronger unpredictability/indistinguishability expectations under extended batteries. Therefore, the current implementation should be treated with caution for cryptographic use until the identified sequence-level issues are addressed and revalidated.


\section{Independent Reproducibility Report of the RCTM-based PRNG}
\label{sec:repro_report}

This section summarizes an independent reproducibility assessment of the Robust Chaotic Tent Map (RCTM) random number generator implementation available in this repository. The objective is to verify whether the observed statistical behavior is consistent across independent execution, and to document both positive and negative findings with full transparency.

\subsection{Repository and Data Availability}

All source code, test scripts, raw test results, and documentation are publicly available at:

\texttt{https://github.com/angualberto/rctm-reproducibility-study}

The repository includes complete test batteries, binary data files (\texttt{agle\_v2\_10m.bin}, \texttt{rctm61.81\_c.bin}), Dieharder results, LaTeX source for this report, and comprehensive documentation (README, quickstart guide, research findings). This enables full independent verification and replication of all findings presented in this study.

\subsection{Experimental Setup}

The experiments were executed on Linux x86\_64 using IEEE-754 double precision arithmetic and C/Python tooling available in the repository. The main analyzed stream was \texttt{agle\_v2\_10m.bin} (10 MB), with additional statistics extracted from:

\begin{itemize}
\item Dieharder full battery (\texttt{-a}): \texttt{agle\_v2\_dieharder\_rctm.txt}
\item Python statistical scripts: \texttt{test\_*.py}
\item Consolidated local report: \texttt{RESULTADOS\_TESTES\_PYTHON\_2026-03-06.txt}
\end{itemize}

\subsection{Technique and Evaluation Protocol Used}

The technical protocol combined first-order distribution checks, sequence-level dependency tests, and distinguishability analysis. The adopted workflow was:

\begin{enumerate}
\item \textbf{Data generation and ingestion}: binary output stream generated by the RCTM implementation and parsed as \texttt{uint8} and \texttt{uint32} depending on each test.
\item \textbf{First-order randomness checks}: Shannon entropy and bit-balance verification to quantify marginal distribution quality.
\item \textbf{Order-pattern complexity}: permutation entropy (order 3) to assess ordinal diversity in local windows.
\item \textbf{Predictability assessment}: N-gram next-symbol prediction (2-gram, 3-gram, 4-gram) with comparison against random baseline (1/256 = 0.391\%).
\item \textbf{Distinguishing capability}: pattern-window entropy test to evaluate whether the generator can be statistically distinguished from an ideal random source.
\item \textbf{Dependency structure}: LAG-1 pair-space coverage and linear autocorrelation across lags 1, 2, 4, 8, 16, and 32.
\item \textbf{State-space probe}: nearest-neighbor state-reconstruction indicator in embedded space.
\item \textbf{External battery validation}: Dieharder \texttt{-a} full suite for broad-spectrum stress testing.
\end{enumerate}

The interpretation criteria followed practical thresholds used in the local scripts: for N-gram, values substantially above baseline indicate exploitable dependence; for autocorrelation, values near zero across multiple lags indicate absence of strong linear dependence; for LAG-1 coverage, low occupancy of the 256$^2$ pair space indicates constrained transitions.

\subsection{Summary of Main Results}

Table~\ref{tab:repro_summary} presents the key outcomes obtained in the independent rerun.

\begin{table}[htbp]
\centering
\caption{Independent reproducibility summary for RCTM output.}
\label{tab:repro_summary}
\begin{tabular}{lcc}
\hline
Test / Indicator & Measured value & Assessment \\
\hline
Shannon entropy & 7.998454 bits/byte & PASS \\
Permutation entropy (order 3) & 0.9999998568 & PASS \\
N-gram predictability (2-gram) & 1.630\% & FAIL \\
Distinguishing attack (normalized entropy) & 0.5190 & FAIL \\
LAG-1 pair coverage & 78.46\% & FAIL \\
Bit balance (ones ratio) & 49.99\% & PASS \\
Autocorrelation (lags 1..32) & near 0 & PASS \\
State reconstruction indicator & no recoverable state & PASS \\
\hline
Overall (Python battery) & 6/9 passed & MIXED \\
\hline
\end{tabular}
\end{table}

In Dieharder, multiple failures/weak results were observed, including \texttt{diehard\_opso}, \texttt{diehard\_squeeze}, repeated failures in \texttt{marsaglia\_tsang\_gcd}, multiple failures in \texttt{rgb\_lagged\_sum}, and failures in \texttt{dab\_bytedistrib}, \texttt{dab\_filltree2}, and \texttt{dab\_monobit2}.

\subsection{Interpretation}

The independent reproduction indicates a \textbf{local-versus-structural randomness gap}. Specifically:

\begin{itemize}
\item \textbf{Local quality appears strong}: high Shannon entropy, good bit balance, and low linear autocorrelation.
\item \textbf{Sequence-level structure remains detectable}: failing N-gram predictability, low pair-space coverage (LAG-1), and distinguishability from ideal random behavior.
\end{itemize}

This combination suggests that, although single-byte statistics are close to ideal, there are non-linear or higher-order dependencies across consecutive outputs.

\subsection{Reproducibility Statement}

All scripts, logs, and result files used to obtain these numbers are included in the repository, enabling direct independent verification. No filtering of unfavorable outcomes was applied in this report.

\subsection{Conclusion for the Reproducibility Study}

From the independent execution, the evidence supports the following conclusion: the evaluated RCTM stream shows good first-order statistical properties but does not consistently satisfy stronger unpredictability/indistinguishability expectations under extended batteries. Therefore, the current implementation should be treated with caution for cryptographic use until the identified sequence-level issues are addressed and revalidated.


\section{Independent Reproducibility Report of the RCTM-based PRNG}
\label{sec:repro_report}

This section summarizes an independent reproducibility assessment of the Robust Chaotic Tent Map (RCTM) random number generator implementation available in this repository. The objective is to verify whether the observed statistical behavior is consistent across independent execution, and to document both positive and negative findings with full transparency.

\subsection{Repository and Data Availability}

All source code, test scripts, raw test results, and documentation are publicly available at:

\texttt{https://github.com/angualberto/rctm-reproducibility-study}

The repository includes complete test batteries, binary data files (\texttt{agle\_v2\_10m.bin}, \texttt{rctm61.81\_c.bin}), Dieharder results, LaTeX source for this report, and comprehensive documentation (README, quickstart guide, research findings). This enables full independent verification and replication of all findings presented in this study.

\subsection{Experimental Setup}

The experiments were executed on Linux x86\_64 using IEEE-754 double precision arithmetic and C/Python tooling available in the repository. The main analyzed stream was \texttt{agle\_v2\_10m.bin} (10 MB), with additional statistics extracted from:

\begin{itemize}
\item Dieharder full battery (\texttt{-a}): \texttt{agle\_v2\_dieharder\_rctm.txt}
\item Python statistical scripts: \texttt{test\_*.py}
\item Consolidated local report: \texttt{RESULTADOS\_TESTES\_PYTHON\_2026-03-06.txt}
\end{itemize}

\subsection{Technique and Evaluation Protocol Used}

The technical protocol combined first-order distribution checks, sequence-level dependency tests, and distinguishability analysis. The adopted workflow was:

\begin{enumerate}
\item \textbf{Data generation and ingestion}: binary output stream generated by the RCTM implementation and parsed as \texttt{uint8} and \texttt{uint32} depending on each test.
\item \textbf{First-order randomness checks}: Shannon entropy and bit-balance verification to quantify marginal distribution quality.
\item \textbf{Order-pattern complexity}: permutation entropy (order 3) to assess ordinal diversity in local windows.
\item \textbf{Predictability assessment}: N-gram next-symbol prediction (2-gram, 3-gram, 4-gram) with comparison against random baseline (1/256 = 0.391\%).
\item \textbf{Distinguishing capability}: pattern-window entropy test to evaluate whether the generator can be statistically distinguished from an ideal random source.
\item \textbf{Dependency structure}: LAG-1 pair-space coverage and linear autocorrelation across lags 1, 2, 4, 8, 16, and 32.
\item \textbf{State-space probe}: nearest-neighbor state-reconstruction indicator in embedded space.
\item \textbf{External battery validation}: Dieharder \texttt{-a} full suite for broad-spectrum stress testing.
\end{enumerate}

The interpretation criteria followed practical thresholds used in the local scripts: for N-gram, values substantially above baseline indicate exploitable dependence; for autocorrelation, values near zero across multiple lags indicate absence of strong linear dependence; for LAG-1 coverage, low occupancy of the 256$^2$ pair space indicates constrained transitions.

\subsection{Summary of Main Results}

Table~\ref{tab:repro_summary} presents the key outcomes obtained in the independent rerun.

\begin{table}[htbp]
\centering
\caption{Independent reproducibility summary for RCTM output.}
\label{tab:repro_summary}
\begin{tabular}{lcc}
\hline
Test / Indicator & Measured value & Assessment \\
\hline
Shannon entropy & 7.998454 bits/byte & PASS \\
Permutation entropy (order 3) & 0.9999998568 & PASS \\
N-gram predictability (2-gram) & 1.630\% & FAIL \\
Distinguishing attack (normalized entropy) & 0.5190 & FAIL \\
LAG-1 pair coverage & 78.46\% & FAIL \\
Bit balance (ones ratio) & 49.99\% & PASS \\
Autocorrelation (lags 1..32) & near 0 & PASS \\
State reconstruction indicator & no recoverable state & PASS \\
\hline
Overall (Python battery) & 6/9 passed & MIXED \\
\hline
\end{tabular}
\end{table}

In Dieharder, multiple failures/weak results were observed, including \texttt{diehard\_opso}, \texttt{diehard\_squeeze}, repeated failures in \texttt{marsaglia\_tsang\_gcd}, multiple failures in \texttt{rgb\_lagged\_sum}, and failures in \texttt{dab\_bytedistrib}, \texttt{dab\_filltree2}, and \texttt{dab\_monobit2}.

\subsection{Interpretation}

The independent reproduction indicates a \textbf{local-versus-structural randomness gap}. Specifically:

\begin{itemize}
\item \textbf{Local quality appears strong}: high Shannon entropy, good bit balance, and low linear autocorrelation.
\item \textbf{Sequence-level structure remains detectable}: failing N-gram predictability, low pair-space coverage (LAG-1), and distinguishability from ideal random behavior.
\end{itemize}

This combination suggests that, although single-byte statistics are close to ideal, there are non-linear or higher-order dependencies across consecutive outputs.

\subsection{Reproducibility Statement}

All scripts, logs, and result files used to obtain these numbers are included in the repository, enabling direct independent verification. No filtering of unfavorable outcomes was applied in this report.

\subsection{Conclusion for the Reproducibility Study}

From the independent execution, the evidence supports the following conclusion: the evaluated RCTM stream shows good first-order statistical properties but does not consistently satisfy stronger unpredictability/indistinguishability expectations under extended batteries. Therefore, the current implementation should be treated with caution for cryptographic use until the identified sequence-level issues are addressed and revalidated.
